\documentclass[11pt]{amsart}

\usepackage[T1]{fontenc}
\usepackage{amsmath,amssymb}
\usepackage{microtype}
\usepackage[
  colorlinks=true,
  linkcolor=blue,
  citecolor=blue,
  urlcolor=blue,
  pdftitle={DRAFT: Minimum Degree Eight and a Rooted-Web Reduction in K7-minus-Minor-Free Contraction-Critical Graphs},
  pdfauthor={Cavit Erginsoy}
]{hyperref}
\setlength{\emergencystretch}{1em}

\newtheorem{theorem}{Theorem}[section]
\newtheorem{lemma}[theorem]{Lemma}
\newtheorem{corollary}[theorem]{Corollary}
\newtheorem{proposition}[theorem]{Proposition}
\newcommand{\minor}{\preccurlyeq}
\newcommand{\N}{N}

\title[DRAFT: Minimum Degree Eight and a Rooted-Web Reduction]
{\textbf{DRAFT}\\[0.5ex]
Minimum Degree Eight and a Rooted-Web Reduction in
$K_7^{-}$-Minor-Free Contraction-Critical Graphs}
\author{Cavit Erginsoy}
\date{Historical draft, 9 August 2026}

\subjclass[2020]{05C15, 05C40}
\keywords{Hadwiger's conjecture, graph minors, contraction-critical graphs, Kempe chains}

\begin{document}

\begin{abstract}
Let $G$ be a minor-minimal non-six-colourable $K_7^{-}$-minor-free graph,
with $n_i$ vertices of degree $i$, and put
$\tau=\sum_{i\geq10}(i-9)n_i$.  We prove
\[
 n_7=0,\qquad \delta(G)\geq8,\qquad |E(G)|\geq4|V(G)|,
 \qquad n_8\geq25+\tau,
\]
and show that $G$ has no $K_5$ subgraph.  Moreover, $G$ has an independent
four-set $U$ of degree-eight vertices such that $G-U$ is three-connected,
nonplanar, and six-chromatic.  For $r\in U$ and a six-colouring of $G-r$,
choose four colours occurring exactly once on $N_G(r)$.  Either their
representatives root a $K_4$ minor in $G-U$, or $G-U$ is a spanning subgraph
of a web rooted at them with a seven-vertex separator
$U\mathbin{\dot\cup}T$.  In the web case, the
three-separations form a lattice after two vertices are fixed on opposite
sides, one of two alternating bichromatic components meets $U-\{r\}$, and
$G[T]\ne K_3$.  Deleting a selected edge yields either five controlled paths
or a smaller component behind another seven-vertex separator, with a boundary
colour class preserved in the latter case.  In either edge-deletion outcome,
there is also a six-ended fan whose residual paths, after removal of its
centre, include at least two pairs with no edge between them.  We also prove that a
six-connected $n$-vertex graph with a degree-$d$ vertex whose neighbourhood
contains $K_4$ and at least $4n+d-13$ edges contains a $K_7^{-}$ minor.  The
proofs are computation-free apart from the cited determination of $R(4,5)$.
\end{abstract}

\maketitle
\markboth{DRAFT: MINIMUM DEGREE EIGHT AND A ROOTED-WEB REDUCTION}
{CAVIT ERGINSOY: DRAFT}

\begin{center}
\small Historical research draft retained for reference; not for submission.
\end{center}

\section{Introduction}

Hadwiger's conjecture states that every graph of chromatic number $t$
contains a $K_t$ minor~\cite{Hadwiger1943}.  It is known for $t\leq6$, with
the last case due to Robertson, Seymour and Thomas
\cite{RobertsonSeymourThomas1993}, and is open for every $t\geq7$.
Write $K_t^{-}$
for the graph obtained from $K_t$ by deleting one edge, and write
$K_7^{\vee}$ for the graph obtained from $K_7$ by deleting two incident
edges.  Every $K_7^{-}$-minor-free graph is
seven-colourable~\cite[Theorem~1]{Albar2014}; Albar's proof uses
Jakobsen's extremal theorem.  Norin and Totschnig proved
that every $K_7^{\vee}$-minor-free graph is
six-colourable~\cite[Theorem~4]{NorinTotschnig2025}.  They asked whether
the same conclusion holds under the weaker exclusion of
$K_7^{-}$~\cite[Conjecture~21]{NorinTotschnig2025}.

Neighbourhoods of low-degree vertices in contraction-critical graphs have
a long history.  Mayer obtained a finite catalogue of candidate
neighbourhood graphs for degree-seven vertices in six-contraction-critical
graphs~\cite{Mayer1993}.  In the minor-minimal
$K_7^{\vee}$-minor-free setting, Norin and Totschnig proved that every
degree-seven vertex belongs to a $K_5$ contained in its closed
neighbourhood~\cite[Claim~4.4]{NorinTotschnig2025}.  Here a rooted-minor
argument first determines the entire neighbourhood.  A second rooted
argument, using their $K^*_{4,2}$ bound, then shows that no
degree-seven vertex can occur at all.

Rolek and Song proved the next theorem for $r=6$
\cite[Lemma~1.9]{RolekSong2017}.  We prove it for every $r\geq2$.

\begin{theorem}\label{thm:two-cliques-intro}
For every $r\geq2$, every $(r+1)$-connected graph containing two distinct
$K_r$ subgraphs contains a $K_{r+2}^{-}$ minor.
\end{theorem}

The connectivity hypothesis is sharp:
$K_{r+1}$ is $r$-connected and contains two distinct $K_r$ subgraphs, but
has no $K_{r+2}^{-}$ minor.

The rooted closure below is the main local tool.  A graph is
\emph{minor-minimal non-six-colourable} if it is not six-colourable and
every proper minor is six-colourable.

\begin{theorem}\label{thm:rooted-helper-intro}
Let $G$ be a six-connected graph of order $n$, let $v\in V(G)$ have degree
$d\geq5$, and suppose that $G[N(v)]$ contains a $K_4$ subgraph.  If
\[
        |E(G)|\geq4n+d-13,
\]
then $K_7^{-}\minor G$.
\end{theorem}

The closure gives the following consequences for a critical graph.

\begin{theorem}\label{thm:package-intro}
Let $G$ be a minor-minimal non-six-colourable $K_7^{-}$-minor-free graph.
For each $i$, let $n_i$ denote the number of vertices of degree $i$, and put
\[
        \tau=\sum_{i\geq10}(i-9)n_i.
\]
Then $G$ is seven-connected and
\[
        n_7=0,
        \qquad \delta(G)\geq8,
        \qquad |E(G)|\geq4|V(G)|.
\]
Moreover, $G$ contains no $K_5$ subgraph and
\[
        n_8\geq25+\tau.
\]
Every degree-eight vertex of $G$ has a $K_4$-free neighbourhood.
\end{theorem}

Four independent degree-eight vertices yield a simultaneous reduction.  A
minor model is \emph{rooted} at nominated vertices when its
branch sets contain them one-to-one.

\begin{theorem}\label{thm:four-centre-intro}
Let $G$ be as in Theorem~\ref{thm:package-intro}.  There is an independent
four-set $U$ of degree-eight vertices.  Put $H=G-U$.  Then $H$ is
three-connected, nonplanar, and six-chromatic, and
$\kappa(G)\in\{7,8\}$.

Fix $r\in U$, a six-colouring $\varphi$ of $G-r$, and four colours which
occur exactly once on $N_G(r)$, with representatives
$x_1,x_2,x_3,x_4$.  Exactly one of the following holds.

\begin{enumerate}
\item The graph $H$ has a $K_4$ minor rooted at
      $x_1,x_2,x_3,x_4$.  Together with the singleton branch set $\{r\}$,
      this is a $K_5$ minor in $G-(U-\{r\})$ with branch sets containing
      $r,x_1,x_2,x_3,x_4$, respectively.
\item The graph $H$ is a spanning subgraph of a web rooted at
      $x_1,x_2,x_3,x_4$.  A facial three-set $T$ of the containing web has
      a nonempty connected set $C$ behind it such that
      \[
             N_H(C)=T,\qquad
             N_G(C)=U\mathbin{\dot\cup}T.              \tag{1.1}
      \]
      Here $T$ is not asserted to be a triangle in $G$; in fact it is not.
\end{enumerate}

If $\kappa(G)=8$, only the first outcome occurs.  If $\kappa(G)=7$, every
proper three-separation of $H$ lifts to an order-seven separation whose
separator consists of $U$ together with its three-vertex separator in $H$.
Oriented between
two fixed vertices lying on opposite open sides of every separation, these
separations form a meet/join lattice.  If three such separations have all eight
sign regions nonempty, each member of $U$ has exactly one neighbour in each
region; four separations cannot have all sixteen sign regions nonempty.  Moreover,
in the web outcome an alternating bichromatic component passes through a
member of $U-\{r\}$.
\end{theorem}

Albar proved the analogous lower bound of twenty-five degree-eight
vertices for a hypothetical minor-minimal non-seven-colourable
$K_7^{-}$-minor-free graph~\cite[Lemma~6]{Albar2014}.  Theorem
\ref{thm:package-intro} concerns criticality one level lower: the rooted
closure eliminates degree seven and gives the defect-sensitive bound
$n_8\geq25+\tau$.

A \emph{minor model} of a graph $H$ in $G$ is a family
$(B_x:x\in V(H))$ of pairwise disjoint connected vertex sets such that
$B_x$ and $B_y$ are adjacent whenever $xy\in E(H)$.  We use
$H\minor G$ to denote that $H$ is a minor of $G$.  All graphs are finite
and simple.

\section{Two linked cliques}

\begin{lemma}\label{lem:disjoint-cliques}
Let $r\geq2$.  If an $(r+1)$-connected graph $H$ contains two
vertex-disjoint $K_r$ subgraphs, then $K_{r+2}^{-}\minor H$.  The branch
sets can be chosen to meet the union of the two cliques.
\end{lemma}

\begin{proof}
Write the two cliques as
\[
        A=\{a_1,\ldots,a_r\},
        \qquad B=\{b_1,\ldots,b_r\}.
\]
By Menger's theorem there are $r$ pairwise vertex-disjoint $A$ to $B$
paths whose open interiors avoid $A\cup B$.  Relabel them as paths $P_i$
from $a_i$ to $b_i$, and put $U=\bigcup_iV(P_i)$.

Call a path $Q$ a \emph{usable cross-path} if its ends lie on distinct
paths $P_i,P_j$, its open interior avoids $U$, and its ends are not both
$A$-ends or both $B$-ends.  Suppose first that $Q$ exists.  Split each of
$P_i$ and $P_j$ across an edge into a nonempty $A$-prefix and a nonempty
$B$-suffix, choosing the split so that the ends of $Q$ belong to pieces of
opposite kinds.  Absorb the open interior of $Q$ into the piece containing
one end.  The four resulting pieces are connected and disjoint.  The two
$A$-pieces are adjacent through $A$, the two $B$-pieces are adjacent
through $B$, the two pieces from each original path are adjacent across
its splitting edge, and $Q$ supplies one of the two remaining
cross-adjacencies.  Thus at most one adjacency among the four pieces is
missing.  Each unsplit path is adjacent to every split piece through its
two clique ends, and the unsplit paths are pairwise adjacent.  These
$r+2$ sets form a $K_{r+2}^{-}$ model.

Suppose now that no usable cross-path exists.  Let $C$ be a component of
$H-U$.  If $C$ has neighbours on two distinct paths, a path through $C$
between two such attachments is unusable only when both attachments are
$A$-ends or both are $B$-ends.  Comparing all pairs of attachments shows
that all neighbours of $C$ are $A$-ends or all are $B$-ends.  Hence
$|\N_H(C)|\leq r$, contrary to $(r+1)$-connectivity.  Every component of
$H-U$ therefore attaches to only one $P_i$.

For each $i$, let $W_i$ be the union of the open interior of $P_i$ and all
components of $H-U$ that attach to $P_i$.  No vertex of $W_i$ has a
neighbour on $P_j$ for $j\ne i$.  Indeed, an edge from the open interior
of $P_i$ to $P_j$ would be a usable cross-path, while such an edge from a
component of $H-U$ would make that component attach to two distinct
linkage paths, contrary to the preceding paragraph.  Hence
\[
        \N_H(W_i)\subseteq\{a_i,b_i\}.
\]
Since $r\geq2$, a nonempty $W_i$ would give a vertex cut of order at most
two.  Thus every $W_i$ is empty.  It follows that $H$ consists of $A$, $B$,
and the matching edges $a_ib_i$: any other edge between $A$ and $B$ would
be usable.  This graph has minimum degree $r$, again contradicting
$(r+1)$-connectivity.
\end{proof}

\begin{proof}[Proof of Theorem~\ref{thm:two-cliques-intro}]
Let $L_1,L_2$ be distinct $K_r$ subgraphs and put
$s=|V(L_1)\cap V(L_2)|$.  If $s\leq r-2$, delete the common set $Z$.
The graph $G-Z$ is $(r+1-s)$-connected and contains the two disjoint
$K_{r-s}$ subgraphs $L_1-Z$ and $L_2-Z$.  Lemma~\ref{lem:disjoint-cliques}
gives a $K_{r-s+2}^{-}$ model in which every branch set contains a vertex
of $(L_1\cup L_2)-Z$.  The vertices of $Z$, taken as singleton branch
sets, are complete to this model and to one another.  They complete a
$K_{r+2}^{-}$ model.

It remains that $s=r-1$.  The union $X=V(L_1)\cup V(L_2)$ induces a graph
containing $K_{r+1}^{-}$.  Since an $(r+1)$-connected graph has at least
$r+2$ vertices, $G-X$ is nonempty.  Let $C$ be a component of $G-X$.  If
$C$ missed a vertex of $X$, then $\N_G(C)$ would be a cut of order at most
$r$.  Thus $C$ is adjacent to every vertex of $X$.  The vertices of $X$
as singleton branch sets, together with $C$, form a $K_{r+2}^{-}$ model.
\end{proof}

We use the case $r=5$ throughout the paper.

\begin{corollary}\label{cor:unique-k5}
Every six-connected $K_7^{-}$-minor-free graph contains at most one $K_5$
subgraph.
\end{corollary}

\section{Degree-seven neighbourhoods}

We use two standard facts about contraction-critical graphs.  A graph is
\emph{$k$-contraction-critical} if it has chromatic number $k$ and every
proper minor is $(k-1)$-colourable.  A noncomplete
$k$-contraction-critical graph is seven-connected when
$k\geq7$~\cite[Satz~6]{Mader1968Connectivity}; see also
\cite[Theorem~1.8]{RolekSong2017}.  Also, if $G$ is $k$-contraction-critical
and $v\in V(G)$, then
\[
        \alpha(G[\N(v)])\leq d(v)-k+2.
                                                        \tag{3.1}
\]
This inequality goes back to Dirac~\cite{Dirac1960}; we use the exact
formulation in \cite[Lemma~1.6(i)]{RolekSong2017}.

We use an immediate specialization of Theorem~7 and property $(*)$ in
Kriesell and
Mohr~\cite[Theorem~7]{KriesellMohr2019}.  Let a graph $J$ be properly
coloured with five colours, and choose one root from each colour class.
Let $D$ be a graph on the five roots with at most six edges.  If the ends
of every edge $xy\in E(D)$ lie in the same bichromatic component of $J$
on their two colours, then $J$ contains pairwise disjoint connected sets,
one containing each root, that are adjacent for every edge of $D$.

\begin{lemma}\label{lem:rooted-k5}
Let $G$ be a seven-chromatic graph such that every proper minor is
six-colourable and $K_7^{-}\not\minor G$.  Let $v$ have degree seven, put
$H=G[\N(v)]$, and let $ab$ be a nonedge of $H$.  If
\[
        R=\N(v)\setminus\{a,b\},
\]
then $G-v-\{a,b\}$ contains a $K_5$ minor rooted at $R$.
\end{lemma}

\begin{proof}
Contract the two edges $va$ and $vb$ to one vertex.  A six-colouring of
the resulting proper minor induces a six-colouring $\varphi$ of $G-v$ in
which $a$ and $b$ have one common colour, say $0$.  Every vertex of $R$
avoids colour $0$.  The five vertices of $R$ use the other five colours
distinctly, since a colour absent from $\N(v)$ could be assigned to $v$.

Put $F=\overline H$.  Inequality~(3.1) gives $\alpha(H)\leq2$, so $F$ is
triangle-free.  If $xy\in E(F[R])$ and $x,y$ have colours $i,j$, then
$x,y$ lie in the same component of the subgraph of $G-v$ induced by
colours $i,j$.  Otherwise, interchanging $i$ and $j$ on the component
containing $x$ would remove colour $i$ from $\N(v)$, allowing that colour
at $v$.

Delete the colour-$0$ class from $G-v$.  The resulting graph is properly
five-coloured, with $R$ a transversal of its colour classes.  By Mantel's
theorem, the triangle-free graph $F[R]$ has at most six edges.  The
Kriesell and Mohr result, applied with demand graph $F[R]$, gives five
disjoint connected sets rooted at $R$ and adjacent for every nonedge of
$H[R]$.  Every other pair is joined by its edge in $H[R]$, so the five
sets form the required rooted $K_5$ model.
\end{proof}

\begin{theorem}\label{thm:local-classification}
Under the hypotheses of Lemma~\ref{lem:rooted-k5}, every degree-seven
vertex $v$ satisfies
\[
        G[\N(v)]\cong K_4\mathbin{\dot\cup}K_3
        \quad\text{or}\quad
        K_1\vee(K_3\mathbin{\dot\cup}K_3).
                                                        \tag{3.2}
\]
In particular, $v$ belongs to a $K_5$ subgraph.
\end{theorem}

\begin{proof}
Put $H=G[\N(v)]$ and $F=\overline H$.  By (3.1), $F$ is triangle-free.
It is nonempty, since otherwise $G[\N[v]]$ would contain $K_7^{-}$.

Every nonisolated vertex of $F$ has degree at least three.  To see this,
choose an edge $ab\in E(F)$ and use the rooted model of
Lemma~\ref{lem:rooted-k5} on $R=\N(v)\setminus\{a,b\}$.  If $d_F(a)=1$,
then $a$ is adjacent in $H$ to every root; the rooted model, $\{a\}$, and
$\{v\}$ form a $K_7$ model.  If $d_F(a)=2$, with other neighbour $x$,
then triangle-freeness gives $bx\in E(H)$.  Absorb $b$ into the branch set
rooted at $x$.  Together with $\{a\}$ and $\{v\}$, the resulting seven
branch sets have at most one missing adjacency, again a contradiction.

Every nontrivial component of $F$ therefore has minimum degree at least
three and hence at least six vertices.  Since $|V(F)|=7$, there is one
nontrivial component and at most one isolated vertex.  If the component
has six vertices, Mantel's theorem and its equality case give
$F\cong K_{3,3}\mathbin{\dot\cup}K_1$.

Suppose instead that the component has seven vertices.  No vertex has
degree at least five: its independent neighbourhood would contain a
vertex of degree at most two.  Thus every degree is three or four.  The
degree sum is even, so some vertex $z$ has degree four.  Its four
neighbours are independent.  Each is adjacent to both vertices outside
$\N_F[z]$ to have degree at least three, and those two outside
vertices are nonadjacent by triangle-freeness.  Hence $F\cong K_{3,4}$.
Taking complements gives (3.2).
\end{proof}

For a minor-minimal non-six-colourable $K_7^{-}$-minor-free graph, Albar's
theorem gives $\chi(G)=7$, while Mader's theorem gives seven-connectivity.
Corollary~\ref{cor:unique-k5}, together with
Theorem~\ref{thm:local-classification}, gives the sharper conclusion below.

\begin{corollary}\label{cor:critical-local}
Let $G$ be as in Theorem~\ref{thm:package-intro}.  Then $G$ contains at
most one $K_5$ subgraph and every degree-seven vertex $v$ satisfies
\[
        G[\N(v)]\cong K_4\mathbin{\dot\cup}K_3.
\]
\end{corollary}

\begin{proof}
The only point not already stated is the exclusion of the second graph in
(3.2).  In that graph, $v$ belongs to two distinct $K_5$ subgraphs, which
contradicts Corollary~\ref{cor:unique-k5}.
\end{proof}

\section[Rooted-model closure]{Rooted
\texorpdfstring{$K^*_{4,2}$}{K-star(4,2)} closure}

For a graph $F$ and $S\subseteq V(F)$, say that $(F,S)$ is
\emph{internally $k$-connected} if there is no separation $(A,B)$ of $F$
with $S\subseteq A$, $B\setminus A\ne\varnothing$, and
$|A\cap B|<k$.  A $Z$-rooted $K^*_{4,2}$ model consists of four pairwise
disjoint connected branch sets, one containing each member of the four-set
$Z$, and two further branch sets $U,V$.  Each of $U,V$ is adjacent to all
four rooted branch sets and to the other; no adjacency between distinct
rooted branch sets is required.

We use the rooted bound of Norin and
Totschnig~\cite[Lemma~12]{NorinTotschnig2025}.

\begin{theorem}[Norin and Totschnig]
\label{thm:rooted-two-helper-bound}
Let $F$ be a graph and let $Z\subseteq V(F)$ have order four.  If $(F,Z)$
is internally four-connected and $F$ has no $Z$-rooted $K^*_{4,2}$ model,
then
\[
        |E(F)|\leq4|V(F)|-10.                            \tag{4.1}
\]
\end{theorem}

The next elementary augmentation places one additional prescribed vertex
in a non-root branch set.

\begin{lemma}\label{lem:fifth-root-helper}
Let $F$ be a graph and let $S=Z\mathbin{\dot\cup}\{x\}$, where $|Z|=4$.
Suppose that $(F,S)$ is internally five-connected.  If $F$ has a
$Z$-rooted $K^*_{4,2}$ model, then it has one in which $x$ belongs to a
non-root branch set.
\end{lemma}

\begin{proof}
Choose a $Z$-rooted model that maximizes $|U|+|V|$ over its two non-root
branch sets $U,V$ and, subject to that, minimizes
$\sum_{i=1}^4|R_i|$ over its rooted branch sets $R_1,\ldots,R_4$.
Write $z_i$ for the root in $R_i$, and put
\[
        P_i=\{r\in R_i:r\text{ has a neighbour in }U\cup V\}.
\]
We claim that $|P_i|=1$.  Both $U$ and $V$ meet $R_i$, so $P_i$ is nonempty.
If it had at least two vertices, there would be distinct $u,w\in R_i$
such that $u$ meets $U$ and $w$ meets $V$; otherwise the two contact sets
would be the same singleton.  Take a minimal tree in $F[R_i]$
containing $z_i,u,w$.  By the minimal choice of the rooted branch sets, this tree
spans $R_i$.  One of $u,w$, say $u$, is a leaf different from $z_i$.
Moving $u$ from $R_i$ into $U$ preserves connectivity of both altered
branch sets.  The former tree edge from $u$ to $R_i\setminus\{u\}$ preserves the
contact between $R_i$ and $U$, while $w$ preserves the contact between
$R_i$ and $V$.  Every
other required adjacency is unchanged, contradicting the maximality of
$|U|+|V|$.  Hence $|P_i|=1$.

No component outside the six branch sets can meet $U\cup V$, since it could
be absorbed into whichever of $U,V$ it meets.  It follows that
\[
        |N_F(U\cup V)\setminus(U\cup V)|\leq4,           \tag{4.2}
\]
with at most one external neighbour in each rooted branch set.  If
$x\notin U\cup V$, (4.2) would give a separation of $(F,S)$ of order at most
four whose second open side is the nonempty set $U\cup V$.  This
contradicts internal five-connectivity.  Thus $x\in U\cup V$.
\end{proof}

\begin{proof}[Proof of Theorem~\ref{thm:rooted-helper-intro}]
Choose a four-set $Z\subseteq N_G(v)$ inducing $K_4$, choose
$x\in N_G(v)\setminus Z$, and put $F=G-v$.  Deleting one vertex from a
six-connected graph leaves a five-connected graph.  In particular,
$(F,Z)$ is internally four-connected and
$(F,Z\cup\{x\})$ is internally five-connected.  Moreover,
\[
 \begin{aligned}
 |E(F)|&=|E(G)|-d\\
       &\geq4n-13
        =4|V(F)|-9.
 \end{aligned}                                           \tag{4.3}
\]
Theorem~\ref{thm:rooted-two-helper-bound} supplies a $Z$-rooted
$K^*_{4,2}$ model in $F$, and Lemma~\ref{lem:fifth-root-helper} lets us
choose it with $x$ in one non-root branch set.

The four rooted branch sets are pairwise adjacent through the edges of
$G[Z]=K_4$.  Together with the two non-root branch sets they form a $K_6$
model.  The singleton branch set $\{v\}$ is adjacent to all four rooted
branch sets and to the one containing $x$; it may miss only the other.
These seven branch sets form a $K_7^{-}$ model in $G$.
\end{proof}

\section{Consequences for a critical graph}

Jakobsen proved that an $n$-vertex
graph with $n\geq7$ and at least $\frac92n-12$ edges contains a
$K_7^{-}$ minor or is a $(K_{2,2,2,2},K_6,4)$-cockade
\cite[Theorem~4]{Jakobsen1983}; see also \cite[Theorem~2]{Albar2014}.

\begin{lemma}\label{lem:jakobsen-defect}
If $G$ is as in Theorem~\ref{thm:package-intro}, with $n=|V(G)|$ and
$m=|E(G)|$, then
\[
        2m\leq9n-25.                                    \tag{5.1}
\]
\end{lemma}

\begin{proof}
The graph $G$ is neither base graph of the cockade family, since both
$K_6$ and $K_{2,2,2,2}$ are six-colourable.  A nontrivial cockade has a
vertex cut of order four, whereas $G$ is seven-connected.  Thus $G$ is not
an exceptional cockade.  Since it has no $K_7^{-}$ minor, Jakobsen's
theorem gives $m<\frac92n-12$.  Integrality yields (5.1).
\end{proof}

The first density bound uses only the linked-clique and neighbourhood
results above.

\begin{proposition}\label{prop:preliminary-density}
Let $G$ be as in Theorem~\ref{thm:package-intro}.  Then
\[
        n_7\leq5,
        \qquad |E(G)|\geq4|V(G)|-2.                     \tag{5.2}
\]
\end{proposition}

\begin{proof}
By Corollary~\ref{cor:critical-local}, every degree-seven vertex lies in
the unique possible $K_5$, so $n_7\leq5$.  Seven-connectivity gives
minimum degree at least seven, and every vertex not counted by $n_7$ has
degree at least eight.  Hence
\[
        2|E(G)|\geq7n_7+8(|V(G)|-n_7)
                  =8|V(G)|-n_7
                  \geq8|V(G)|-5.
\]
The left side is even, so it is at least $8|V(G)|-4$, proving (5.2).
\end{proof}

\begin{proof}[Proof of Theorem~\ref{thm:package-intro}]
Albar's theorem gives $\chi(G)=7$, and every proper minor is
six-colourable by the choice of $G$.  The graph is noncomplete because it
has no $K_7^{-}$ minor, so Mader's theorem gives seven-connectivity.

Suppose that $v$ has degree seven.  Corollary~\ref{cor:critical-local}
shows that $G[N(v)]$ contains a $K_4$.  Proposition
\ref{prop:preliminary-density} gives
\[
        |E(G)|\geq4|V(G)|-2>4|V(G)|-6,
\]
which is more than the threshold in Theorem
\ref{thm:rooted-helper-intro} when $d_G(v)=7$.  That theorem gives the
forbidden $K_7^{-}$ minor.  Thus $n_7=0$.  Seven-connectivity now gives
$\delta(G)\geq8$, and degree summation gives
$|E(G)|\geq4|V(G)|$.

Put $q=|E(G)|-4|V(G)|$, so $q\geq0$.  Suppose that a $K_5$ subgraph has
vertex set $K$.  For each $w\in K$, the neighbourhood of $w$ contains the
$K_4$ on $K\setminus\{w\}$.  Since $G$ has no $K_7^{-}$ minor, the
contrapositive of Theorem~\ref{thm:rooted-helper-intro} gives
\[
        4|V(G)|+q<4|V(G)|+d_G(w)-13,
\]
and hence $d_G(w)\geq q+14$.  The five vertices of $K$ would contribute at
least $5(q+6)$ to
\[
        \sum_{z\in V(G)}(d_G(z)-8)=2q,
\]
which is impossible.  Thus $G$ contains no $K_5$ subgraph.

Finally, Lemma~\ref{lem:jakobsen-defect} and the exact degree identity give
\[
        25\leq9|V(G)|-2|E(G)|=n_8-\tau.
\]
Therefore $n_8\geq25+\tau$.  If a degree-eight vertex had a $K_4$ in its
neighbourhood, it would belong to a $K_5$, which has just been excluded.
\end{proof}

\section{Four independent degree-eight vertices}

We prove Theorem~\ref{thm:four-centre-intro} using $R(5,4)=25$, which
follows by complement symmetry from the theorem of McKay and
Radziszowski~\cite{McKayRadziszowski1995}, and the rooted $K_4$ theorem of
Fabila-Monroy and Wood~\cite{FabilaMonroyWood2013}.

\begin{lemma}\label{lem:deleted-four-centres}
Let $G$ be as in Theorem~\ref{thm:package-intro}.  Then $G$ contains an
independent four-set $U$ of degree-eight
vertices.  For every such set, $H=G-U$ is three-connected, nonplanar, and
six-chromatic.  Moreover $\kappa(G)\in\{7,8\}$, and if $\kappa(G)=8$ then
$H$ is four-connected.
\end{lemma}

\begin{proof}
By Theorem~\ref{thm:package-intro}, at least twenty-five vertices have
degree eight and $G$ contains no $K_5$ subgraph.  The Ramsey equality
$R(5,4)=25$ therefore supplies the required independent set $U$.

Deleting four vertices lowers connectivity by at most four, so $H$ is
three-connected, and is four-connected when $G$ is eight-connected.  For
each $u\in U$, the set $N_G(u)$ separates $u$ from the other members of
$U$, whence $\kappa(G)\leq8$.

The independence and degrees of the four vertices give
\[
 |E(H)|=|E(G)|-32\geq4|V(H)|-16.
\]
Since $|V(H)|\geq21$, this exceeds $3|V(H)|-6$, so $H$ is nonplanar.
It is six-colourable as a proper minor of $G$.  A five-colouring of $H$
would extend to a six-colouring of $G$ by assigning one new colour to all
of $U$.  Thus $\chi(H)=6$.
\end{proof}

For the rest of this section, let $G$ be as in
Theorem~\ref{thm:package-intro}, let $U$ be an independent four-set of
degree-eight vertices, and put $H=G-U$.  Statements involving
$r,\varphi,x_1,x_2,x_3,x_4$ use the choices in
Theorem~\ref{thm:four-centre-web}; in its web outcome, $T,C$, and
$S=U\mathbin{\dot\cup}T$ have the meanings fixed there.

\subsection{The rooted-web alternative}

A web
rooted at four nominated vertices has a plane skeleton with the roots on
its outer face; additional clique sets may be placed behind facial
triangles and joined to all three vertices of the corresponding triangle.
We call the three vertices of such a triangle a \emph{facial three-set};
they need not induce a triangle in a spanning subgraph of the web.
Fabila-Monroy and Wood proved that a three-connected graph has a $K_4$
minor rooted at four nominated vertices if and only if it is not a spanning
subgraph of a web rooted at those vertices
\cite[Theorem~8]{FabilaMonroyWood2013}.  Edges of a facial triangle in the
containing web need not be edges of the spanning subgraph.

\begin{theorem}\label{thm:four-centre-web}
Fix $r\in U$ and a proper six-colouring $\varphi$ of $G-r$.  Choose four
colours occurring exactly once on $N_G(r)$, and let
$x_1,x_2,x_3,x_4$ be their unique representatives.  Exactly one of the
following holds.

\begin{enumerate}
\item $H$ contains a $K_4$ model rooted at $x_1,x_2,x_3,x_4$.
      Together with $\{r\}$, this is a $K_5$ model in
      $G-(U-\{r\})$.
\item $H$ is a spanning subgraph of a web rooted at
      $x_1,x_2,x_3,x_4$.  Some facial three-set $T$ and some nonempty
      connected set $C\subseteq V(H)-T$ satisfy
      \[
             N_H(C)=T,\qquad
             N_G(C)=S:=U\mathbin{\dot\cup}T.           \tag{6.1}
      \]
      No $x_i$ belongs to $C$, and at least one $x_i$ belongs to a component
      of $G-S$ other than $C$.  Every component $D$ of $G-S$ satisfies
      $N_G(D)=S$.
\end{enumerate}
\end{theorem}

\begin{proof}
Every colour occurs on $N_G(r)$, since otherwise $\varphi$ extends to
$r$.  Six positive multiplicities summing to eight include at least four
ones.  Since $U$ is independent, the corresponding representatives lie in
$H$.

Apply the rooted theorem quoted above to $H$.  Its rooted-model outcome
gives item 1 because $r$ is adjacent to the nominated vertex in every rooted
branch set.  All five branch sets avoid $U-\{r\}$.

Otherwise $H$ is a spanning subgraph of the stated web.  Since $H$ is
nonplanar, some clique set $X_T$ behind a facial three-set $T$ is nonempty
in $H$.  Let $C$ be a component of $H[X_T]$.  Web containment gives
$N_H(C)\subseteq T$.
At least one outer root lies outside $T$, so three-connectivity forces
$N_H(C)=T$.

In $G$, every neighbour of $C$ lies in $T\cup U$.  The outer root chosen
outside $T$ also lies outside $C\cup T\cup U$, so $N_G(C)$ separates it
from $C$.  Seven-connectivity therefore forces $N_G(C)=T\cup U$, proving
(6.1).  The vertices added behind $T$ are disjoint from the skeleton, so no
$x_i$ belongs to $C$, and one lies outside the three-set $T$.  Finally, any
component of $G-S$ whose
neighbourhood were a proper subset of $S$ would contradict
seven-connectivity.  Thus every component $D$ of $G-S$ has $N_G(D)=S$.
\end{proof}

\begin{corollary}\label{cor:eight-connected-four-centre}
If $\kappa(G)=8$, item 1 of Theorem~\ref{thm:four-centre-web} holds for
every choice of $r$, $\varphi$, and the four singleton-colour neighbours.
\end{corollary}

\begin{proof}
Item 2 gives a seven-vertex separator.
\end{proof}

\subsection{Minimum lift order}

For $u\in U$, write $N_H(u)=N_G(u)\cap V(H)$.  For a separation
$p=(A,B)$ of $H$, write
\[
 L_p=A-B,\qquad R_p=B-A,\qquad S_p=A\cap B,
\]
and say that $u$ \emph{crosses} $p$ if it has a neighbour in both $L_p$
and $R_p$.  Define
\begin{align}
 C_U(p)&=\{u\in U:N_H(u)\cap L_p\ne\varnothing\ne
                         N_H(u)\cap R_p\},\notag\\
 \lambda_U(p)&=|S_p|+|C_U(p)|.                         \tag{6.2}
\end{align}
A lift of $p$ is \emph{trace-preserving} if its intersections with $H$ are
exactly $A$ and $B$.

\begin{theorem}\label{thm:lift-order}
For every separation $p$ of $H$, the number $\lambda_U(p)$ is the minimum
order of a trace-preserving lift to $G$.  The function $\lambda_U$ is
symmetric and submodular.  More precisely, for separations $p=(A,B)$ and
$q=(C,D)$, put
\[
 p\wedge q=(A\cap C,B\cup D),\qquad
 p\vee q=(A\cup C,B\cap D).
\]
Then
\[
 \lambda_U(p\wedge q)+\lambda_U(p\vee q)
       \leq\lambda_U(p)+\lambda_U(q).                  \tag{6.3}
\]
If equality holds, equality holds separately for the crossing indicator
of every $u\in U$.  If $G$ is $k$-connected and $p$ is proper, then
$\lambda_U(p)\geq k$.
\end{theorem}

\begin{proof}
Every lift must put $S_p$ and every member of $C_U(p)$ in its separator.
Conversely, put precisely those vertices there; put each remaining member
of $U$ on a side containing all its neighbours outside $S_p$.
The independence of $U$ creates no crossing edge.  This proves the minimum
formula and also symmetry.

The ordinary separator terms satisfy the modular identity
\[
 |S_p|+|S_q|=|S_{p\wedge q}|+|S_{p\vee q}|.            \tag{6.4}
\]
For one fixed $u$, its crossing indicator is submodular.  Indeed,
\[
\begin{aligned}
 L_{p\wedge q}&=L_p\cap L_q,&R_{p\wedge q}&=R_p\cup R_q,\\
 L_{p\vee q}&=L_p\cup L_q,&R_{p\vee q}&=R_p\cap R_q.
\end{aligned}
\]
If $u$ crosses both corners it crosses both inputs; if it crosses only one
corner it crosses at least one input.  Summing over $U$ and using (6.4)
proves (6.3), and equality in the sum implies equality for each $u\in U$.
A proper separation of $H$ has a proper minimum lift, so $k$-connectivity
gives the last assertion.
\end{proof}

For a family of oriented separations, fix vertices $a,b$ such that
$a\in L_p$ and $b\in R_p$ for every separation $p$ in the family.

\begin{corollary}\label{cor:four-centre-lattice}
Assume $\kappa(G)=7$.  Every proper three-separation $p$ of $H$ is crossed
by all four members of $U$ and lifts to an order-seven separation with
separator $U\mathbin{\dot\cup}S_p$, where $S_p$ is its three-vertex
separator in $H$.  Conversely, the trace of every order-seven separation of
this form is a proper three-separation of $H$ crossed by all four members of
$U$.

For fixed $a,b$ as above, these oriented separations form a finite meet/join
lattice.  Its inclusion-minimum open side containing $a$ is connected.
\end{corollary}

\begin{proof}
For a proper three-separation, Theorem~\ref{thm:lift-order} gives
$3+|C_U(p)|\geq7$, so all four members of $U$ cross and equality holds.
Conversely, every vertex in the separator of an order-seven separation of a
seven-connected graph has a neighbour in both open sides; otherwise the
other six separator vertices would still separate one side.

If two order-seven separations have the same fixed vertices $a,b$, their
meet and join are proper.  Connectivity bounds each corner order below by
seven, while submodularity bounds their sum above by fourteen.  Both corners
therefore have order seven.  Equality in (6.3), together with the pointwise
equality clause of Theorem~\ref{thm:lift-order}, makes every member of $U$
cross both corners.  Hence each corner has $H$-separator order $7-4=3$, so
repeated meets stay in the stated family and give the minimum open side.  If
that side were disconnected, the component containing $a$ would have
neighbourhood contained in the same seven-set.  Seven-connectivity makes
that neighbourhood the whole separator, producing a smaller open side
containing $a$, a contradiction.
\end{proof}

\begin{corollary}\label{cor:crossing-regions}
Let $p_1,\ldots,p_m$ be proper three-separations of $H$ whose
trace-preserving lifts have order seven as in
Corollary~\ref{cor:four-centre-lattice}.  For
$\sigma\in\{+,-\}^m$, let $R_\sigma$ be the intersection of the open
sides selected by its signs.  If both $R_\sigma$ and $R_{-\sigma}$ are
nonempty, then every $u\in U$ has a neighbour in $R_\sigma$.

Consequently at most eight sign regions whose antipodal regions are also
nonempty can occur.  If all four regions of two separations are nonempty,
every member of $U$ meets all four.  If all eight regions of three
separations are nonempty, every member of $U$ has exactly one neighbour in
each region and no separator neighbour.  Four separations cannot have all
sixteen sign regions nonempty.
\end{corollary}

\begin{proof}
Choose $a\in R_\sigma$ and $b\in R_{-\sigma}$, orient the separations
accordingly, and take repeated meets.  Corollary
\ref{cor:four-centre-lattice} gives another order-seven separation whose
selected open side is $R_\sigma$; every member of $U$ crosses it.  Distinct
regions are disjoint and every member of $U$ has degree eight, giving all
the conclusions.
\end{proof}

\subsection{The colour-labelled web obstruction}

\begin{theorem}\label{thm:centre-kempe}
Assume the web outcome and list $x_1,x_2,x_3,x_4$ in their cyclic order on
the outer face.  Put $c_i=\varphi(x_i)$.  Every pair $x_i,x_j$ lies in the
same $(c_i,c_j)$-component of $G-r$.  Let $K_{13}$ be the
$(c_1,c_3)$-component
containing $x_1,x_3$, and define $K_{24}$ analogously.  Then $K_{13}$ and
$K_{24}$ are vertex-disjoint, and at least one contains a member of
$U-\{r\}$.

If $\Delta$ is the set of two unselected colours, some
$s\in U-\{r\}$ sees both colours of $\Delta$ in $H$.
\end{theorem}

\begin{proof}
If $x_i,x_j$ were in different $(c_i,c_j)$-components, interchange the two
colours on the component containing $x_i$.  Since $x_i$ and $x_j$ are the
unique neighbours of $r$ in their respective colours, colour $c_i$ would
then be absent from $N_G(r)$, and could be assigned to $r$.  This is
impossible.

The two components $K_{13}$ and $K_{24}$ use disjoint colour pairs and
hence are vertex-disjoint.  If both avoided $U-\{r\}$, they would give in
$H$ two disjoint paths joining alternating outer pairs, contrary to the
Two Paths description of a web
\cite[Lemma~2 and the first proof of Lemma~4]{FabilaMonroyWood2013}.

Suppose finally that every $s\in U-\{r\}$ missed in $N_H(s)$ some colour
$\delta_s\in\Delta$.  Recolour each $s$ with $\delta_s$.  These changes
are simultaneous and proper because $U$ is independent, and they do not
change the colours on $N_G(r)$.  Repeating the preceding argument, both
alternating components now avoid $U-\{r\}$, the same contradiction.
\end{proof}

\subsection{The facial three-set}

\begin{corollary}\label{cor:web-boundary-residue}
In the web outcome, $G[T]$ is not a triangle.
\end{corollary}

\begin{proof}
Suppose that $G[T]$ is a triangle, say $T=\{t_1,t_2,t_3\}$.  Let
$R=V(G)-(C\cup S)$.  Since $G$ is seven-connected, every component of
$G-S$ has neighbourhood $S$, so both $C\cup U$ and $R\cup U$ are connected.
Contract a spanning tree of $C\cup U$, six-colour the proper minor, and pull
the colouring back only to $G[R\cup S]$.  Its colour classes on $S$ are
\[
                         U\mid\{t_1\}\mid\{t_2\}
                           \mid\{t_3\}.                \tag{6.5}
\]
Contracting $R\cup U$ instead produces the same colour classes on
$G[C\cup S]$.  A permutation of colour names aligns the two colourings,
which then glue to a six-colouring of $G$, a contradiction.
\end{proof}

\subsection{A colouring response to an edge deletion at the web separator}

\begin{theorem}[Controlled paths or a smaller component]
\label{thm:web-edge-response}
Assume outcome 2 of Theorem~\ref{thm:four-centre-web}.  Choose $x=x_j$ in
a component $D$ of $G-S$ distinct from $C$, put
$\gamma=\varphi(x)$, and let $e=rx$.  Thus
$S=U\mathbin{\dot\cup}T$ and $N_G(C)=N_G(D)=S$.
Extend $\varphi$ to $r$ by assigning colour $\gamma$ to $r$, and call the
resulting map $d$.  Then $d$ is a proper six-colouring of $G-e$.

For every colour $\beta\ne\gamma$, the vertices $x$ and $r$ lie in the same
$(\gamma,\beta)$-component of $G-e$.  Choose a simple path from $x$ to $r$
in this component, orient it from $x$, and stop it at its first vertex
$q_\beta\in S$.  Write $xy_\beta$ for its first edge; the stopped path has
all internal vertices in $D$.  Put
\[
 Q_0=\{r\}\cup\{q_\beta:\beta\ne\gamma\}.
\]
The five vertices $y_\beta$ are distinct.  At least one of the following
holds.

\begin{enumerate}
\item There is a set $Q\subseteq S$ with $Q_0\subseteq Q$ and $|Q|\leq6$,
      and five paths from $x$ to $Q$ which begin with the prescribed edges
      $xy_\beta$ and are pairwise vertex-disjoint outside
      $\{x\}\cup Q$.  Their ends in $Q$ need not be distinct.
\item One has $|Q_0|\leq5$, and there are a five-set $Q\subseteq S$
      containing $Q_0$, a four-set $Z\subseteq D-\{x\}$, and a nonempty
      connected set $A\subseteq D-(\{x\}\cup Z)$ such that
      \[
        N_G(A)=(S-Q)\mathbin{\dot\cup}\{x\}
                         \mathbin{\dot\cup}Z.            \tag{6.6}
      \]
      In particular, $N_G(A)$ is a seven-vertex separator and
      $|A|<|D|$.  The restriction of $d$ to
      $G[A\cup N_G(A)]$ is a proper six-colouring.  If
      \[
          J=\{z\in N_G(A):d(z)=\gamma\},
      \]
      then $|J|\geq2$.  In this colouring, and also in some proper
      six-colouring of $G-A$, the set $J$ is exactly one colour class on
      $N_G(A)$.
\end{enumerate}

Independently of which alternative occurs, there is a fan from $x$ to six
distinct vertices of $S$, consisting of the edge $xr$ and five paths
beginning with the prescribed edges $xy_\beta$.  Delete $x$ from the five
non-direct paths and denote the resulting disjoint connected sets by
$L_\beta$.  The contact graph of these five sets, in which two are
adjacent when they have an edge between them in $G$, has at most eight
edges.
\end{theorem}

\begin{proof}
The vertex $x$ is the unique $\gamma$-coloured neighbour of $r$.
Consequently $rx$ is the only monochromatic edge after restoring $r$, so
$d$ is proper on $G-e$.

Fix $\beta\ne\gamma$.  If $x$ and $r$ belonged to different
$(\gamma,\beta)$-components, interchange the two colours on the component
containing $x$.  Every $\beta$-coloured neighbour of $r$ lies in the
component containing $r$, and $x$ is its only $\gamma$-coloured neighbour.
The interchange would therefore repair $rx$ without creating another
conflict at $r$, giving a six-colouring of $G$.  This contradiction proves
the claimed connection.  Since $d(x)=\gamma$, each $y_\beta$ has colour
$\beta$; hence the five vertices $y_\beta$, and the five prescribed edges,
are distinct.

Let $I=\{\beta:y_\beta=q_\beta\}$ be the colours whose prescribed path is
already a boundary edge, and put $q=|I|$.

Suppose first that $|Q_0|\leq5$, and enlarge $Q_0$ to a five-set
$Q\subseteq S$.  Retain the $q$ direct paths.  In
\[
                    G[(D-\{x\})\cup Q]
\]
seek paths from the other $5-q$ vertices $y_\beta$ to $Q$, using every
source, with unit vertex capacity in $D-\{x\}$ and unrestricted capacity
at $Q$.  If such a routing exists, prepending the prescribed first edges
gives outcome 1.

Otherwise, the capacitated form of Menger's theorem gives a set
$Z\subseteq D-\{x\}$ with $|Z|\leq4-q$ meeting every such path.  Some
remaining source avoids $Z$; let $A$ be its component in
$G[D-(\{x\}\cup Z)]$.  It has no neighbour in $Q$, and therefore
\[
             N_G(A)\subseteq(S-Q)\cup\{x\}\cup Z.
\]
The right side has order at most $7-q$.  The nonempty component $C$ lies
on the other side, so this is a genuine separation.  Seven-connectivity
forces $q=0$, $|Z|=4$, and equality throughout, proving (6.6).

For each $\beta\ne\gamma$, the original stopped path with $x$ deleted
joins $y_\beta$ to $Q$, and hence meets $Z$.  Choose one such intersection
vertex $z_\beta$.  Five choices in the four-set $Z$ give
$z_\beta=z_{\beta'}$ for two distinct colours $\beta,\beta'$.  A common
vertex of a $\gamma,\beta$ path and a $\gamma,\beta'$ path has colour
$\gamma$.  Thus $J$ contains both this vertex and $x$, so $|J|\geq2$.

The only edge on which $d$ is improper in $G$ is $rx$.  Since $r\in Q$,
this edge is absent from $G[A\cup N_G(A)]$, so the restriction of $d$ to
that graph is proper.  The set $A\cup J$ is connected, and $J$ is independent.
Contract a spanning tree of $G[A\cup J]$ and six-colour the resulting
proper minor.  Pull the colouring back to $G-A$, assigning every member of
$J$ the colour of the contraction image.  That image is adjacent to every
vertex of $N_G(A)-J$, so $J$ is exactly one colour class on $N_G(A)$.

It remains to treat $|Q_0|=6$.  Set $Q=Q_0$ and retain the $q$ direct
paths as above.  Failure to route the remaining sources to $Q$ would give
a set $Z$ of order at most $4-q$ and a nonempty connected set
$A\subseteq D-\{x\}$ satisfying
\[
             N_G(A)\subseteq(S-Q)\cup\{x\}\cup Z.
\]
This set has order at most $1+1+(4-q)\leq6$, again separates $A$ from
$C$, and contradicts seven-connectivity.  Outcome 1 therefore holds.

For the distinct-ended fan, again retain the $q$ prescribed edges which
already end in $S$.  Put
\[
        W=S-\bigl(\{r\}\cup\{y_\beta:\beta\in I\}\bigr).
\]
Working in $G[(D-\{x\})\cup W]$, route the remaining $5-q$ vertices
$y_\beta$ to distinct members of $W$.  If this routing failed, Menger's theorem would
give a separator of order at most $4-q$.  Since there are $5-q$ sources
and $6-q$ targets, at least one of each survives.  A surviving source
component would then have its
neighbourhood contained in that separator together with $x$, $r$, and the
$q$ retained boundary ends, a set of order at most six.  A surviving
target lies on the other side, contradicting seven-connectivity.
Prepending the prescribed first edges and adjoining $xr$ gives the fan.

Finally suppose that the contact graph of the five sets $L_\beta$ had at
least nine edges.  The seven disjoint connected branch sets
\[
             \{x\},\qquad C\cup\{r\},\qquad
             L_\beta\quad(\beta\ne\gamma)
\]
would form a $K_7^{-}$ model.  The first set meets every $L_\beta$ through
its prescribed first edge.  The set $C\cup\{r\}$ is connected, meets
$\{x\}$ through $rx$, and meets every $L_\beta$ through its distinct end
in $S-\{r\}$ because $N_G(C)=S$.  Among the five remaining
branch sets at most one adjacency is missing.  This contradicts the choice
of $G$, and proves the contact bound.
\end{proof}

\begin{corollary}[Restoring one vertex of $U$]\label{cor:named-centre-repair}
In the web outcome, for every $s\in U-\{r\}$ the graph
$G-(U-\{s\})$ contains a $K_4$ model rooted at
$x_1,x_2,x_3,x_4$, and every such model uses $s$.
\end{corollary}

\begin{proof}
Deleting the three vertices of $U-\{s\}$ from the seven-connected graph
$G$ leaves the four-connected graph $G-(U-\{s\})$.  It is nonplanar because
it contains $H$.  The four-connected rooted-$K_4$ theorem of
Fabila-Monroy and Wood~\cite[Theorem~6]{FabilaMonroyWood2013} therefore
supplies the rooted model.  The graph $H$ has no such model in the web
outcome, so every model in $G-(U-\{s\})$ must use $s$.
\end{proof}

The web case remains open: the smaller separator in
Theorem~\ref{thm:web-edge-response} need not contain $U$, and the fan paths
are not assigned to prescribed branch sets.  One must retain $U$ when passing
to the smaller component or absorb the paths into one rooted model.

\section{Extremal reduction and scope}

\begin{corollary}[Extremal reduction]\label{cor:extremal-reduction}
Suppose that every seven-connected graph $H$ satisfying
\[
        |E(H)|\geq4|V(H)|
\]
contains a $K_7^{-}$ minor.  Then every $K_7^{-}$-minor-free graph is
six-colourable.
\end{corollary}

\begin{proof}
Otherwise, choose a minor-minimal non-six-colourable
$K_7^{-}$-minor-free graph $G$.  By Theorem~\ref{thm:package-intro}, the
graph $G$ is seven-connected and $|E(G)|\geq4|V(G)|$.  The assumed
extremal statement therefore gives $K_7^{-}\minor G$, a contradiction.
\end{proof}

The extremal hypothesis remains open, so this paper proves neither the Norin
and Totschnig conjecture nor Hadwiger's conjecture for $t=7$.

\section*{Statement on AI use and independence}

GPT-5.6 Sol, an OpenAI model, contributed to theorem discovery, proofs,
verification, and drafting.  Cavit Erginsoy directed the work and made the
final decisions.  Erginsoy is not affiliated with OpenAI, and OpenAI neither
sponsored nor verified the work.

{\hbadness=2000
\bibliographystyle{amsplain}
\bibliography{references}
}

\end{document}
